% ============================================================
% CAPITOLO 41 - PROGRESSIONI CROMATICHE
% ============================================================
\chapter{Progressioni cromatiche}
\label{chap:progressioni_cromatiche}

Le progressioni cromatiche introducono note estranee alla scala di partenza. Sono efficaci per creare movimento per semitoni verso l'alto o verso il basso, e i compositori barocchi le hanno usate per costruire tensione o per modulare.

\section{Progressione cromatica ascendente (5, 6/5)}
\label{subsec:progressione_cromatica_ascendente}

Il basso sale per semitoni, e gli accordi alternano triadi in stato fondamentale e settime al primo rivolto.

\begin{regola}[Progressione cromatica ascendente (numerica 5, 6/5)]
	\label{reg:progressione_cromatica_asc}
	\begin{itemize}
		\item Il basso sale per \textbf{semitoni}.
		\item Gli accordi alternano \textbf{triadi in stato fondamentale} (cifra 5) e \textbf{settime al primo rivolto} (cifra 6/5): 5 - 6/5 - 5 - 6/5 - 5.
		\item Le alterazioni cromatiche creano sensibili «artificiali» che tendono a salire.
		\item Costruisce tensione verso un culmine.
	\end{itemize}
	Ogni settima al primo rivolto (6/5) contiene una dissonanza che risolve sull'accordo seguente; l'alterazione crea una «piccola dominante» che spinge il discorso in alto.
\end{regola}

\begin{esempio}[Progressione cromatica ascendente in Do minore]
	\label{ex:progressione_cromatica_asc}
	Il basso sale La$\flat$--La--Si$\flat$--Si--Do. Gli accordi sono VI, \sigla{V}{6}{5} di VII (Fa con settima, sul La), VII, \sigla{V}{6}{5} (Sol con settima, sul Si), i. Il La e il Si naturali creano la dominante secondaria del VII e la dominante, che spingono verso la tonica.
	\musica{es_cap17_13}
\end{esempio}

Perché un cromatismo ascendente accumula tensione è detto nell'Approfondimento \ref{phys:energia_ascendenti}.

\section{Progressione cromatica discendente (2, 7)}
\label{subsec:progressione_cromatica_discendente}

La progressione cromatica discendente è speculare: il basso scende per semitoni e gli accordi di settima si alternano nel terzo rivolto e in stato fondamentale.

\begin{regola}[Progressione cromatica discendente (numerica 2, 7)]
	\label{reg:progressione_cromatica_disc}
	\begin{itemize}
		\item Il basso scende per \textbf{semitoni}.
		\item Gli accordi alternano \textbf{settime al terzo rivolto} (cifra 2) e \textbf{settime in stato fondamentale} (cifra 7): 2 - 7 - 2 - 7 - 2.
		\item Le alterazioni cromatiche (spesso $\flat$) facilitano la discesa.
		\item L'effetto è di una discesa controllata, spesso malinconica o patetica.
	\end{itemize}
	Nel terzo rivolto il basso è la settima dell'accordo, che deve scendere di grado: questo produce il movimento discendente. La successione 2-7 lo mantiene cromatico nel basso, mentre le altre voci si muovono più liberamente.
\end{regola}

In Re minore lo schema di basso più noto scende dalla tonica alla dominante, Re -- Do$\sharp$ -- Do -- Si -- Si$\flat$ -- La: a ogni semitono la settima del terzo rivolto scende, e l'accordo seguente in stato fondamentale la prende come terza o come quinta. La Sezione \ref{sec:basso_cromatico_disc} mostra due armonizzazioni complete del basso cromatico discendente.

Un esempio celebre di basso cromatico discendente è il \textit{Crucifixus} della Messa in Si minore di J.S. Bach, costruito su un basso ostinato. Nella retorica musicale barocca il movimento cromatico discendente è associato al \textit{passus duriusculus} (passo difficile), figura che rappresenta la sofferenza.

Il cromatismo pone una questione legata al temperamento equabile (Approfondimento \ref{math:cromatismo_temperamento}).

\section{Basso cromatico discendente}
\label{sec:basso_cromatico_disc}

Il basso scende di semitono a ogni accordo: Do -- Si -- Si$\flat$ -- La -- La$\flat$ -- Sol. Gli accordi sono in stato fondamentale o rivoltati, in modo che il basso proceda per grado cromatico. Le due armonizzazioni usano solo triadi (I -- \sigla{V}{6}{} -- $\flat$VII -- \sigla{IV}{6}{} -- \sigla{iv}{6}{} -- V) oppure triadi e settime (il terzo accordo diventa \sigla{I}{2}{}, cioè \sigla{V}{2}{}/IV, con la settima Si$\flat$ al basso).

\begin{regola}[Armonizzazione legata e sincopata]
	\label{reg:cromatico_legato_sincopato}
	\begin{itemize}
		\item \textbf{Legato:} le voci superiori procedono con il minimo movimento, tenendo le note comuni (legature).
		\item \textbf{Sincopato:} ogni nota del basso entra sul tempo \emph{debole} e resta legata, oltre la stanghetta, fino al tempo forte successivo; le voci superiori cambiano accordo solo sul tempo forte. Sul tempo debole il basso anticipa quindi la nota dell'accordo seguente.
	\end{itemize}
	Con questa armonizzazione la versione con sole settime non è realizzabile: la settima di Fa minore non si può preparare e la settima maggiore di Do non trova la nota di risoluzione.
\end{regola}

\begin{esempio}[Basso cromatico discendente in Do maggiore, legato]
	\label{ex:cromatico_legato}
	Solo triadi:
	\musica{es_cap17_crom_leg_T}
	Triadi e settime:
	\musica{es_cap17_crom_leg_TS}
\end{esempio}

\begin{esempio}[Basso cromatico discendente in Do maggiore, sincopato]
	\label{ex:cromatico_sincopato}
	Solo triadi:
	\musica{es_cap17_crom_sin_T}
	Triadi e settime:
	\musica{es_cap17_crom_sin_TS}
\end{esempio}
